% Figure 06 — Diagramme de séquence : Authentification
% Acteurs : User, Browser, Nginx, Flask, PostgreSQL
\documentclass[border=10pt]{standalone}
% Préambule commun pour tous les diagrammes TikZ standalone
% À inclure via % Préambule commun pour tous les diagrammes TikZ standalone
% À inclure via % Préambule commun pour tous les diagrammes TikZ standalone
% À inclure via \input{preambule.tex} après \documentclass{standalone}
\usepackage[utf8]{inputenc}
\usepackage[T1]{fontenc}
\usepackage[french]{babel}
\usepackage{lmodern}
\usepackage{tikz}
\usepackage{xcolor}

\usetikzlibrary{
  arrows.meta,
  positioning,
  shapes.geometric,
  shapes.symbols,
  shapes.misc,
  fit,
  calc,
  backgrounds,
  decorations.pathreplacing,
  decorations.markings,
  shadows.blur,
  patterns
}

% Palette de couleurs (cohérente sur tout le rapport)
\definecolor{awsorange}{HTML}{FF9900}
\definecolor{awsdark}{HTML}{232F3E}
\definecolor{awsblue}{HTML}{1A476F}
\definecolor{dockerblue}{HTML}{2496ED}
\definecolor{dockerdark}{HTML}{0DB7ED}
\definecolor{nginxgreen}{HTML}{009639}
\definecolor{flaskgray}{HTML}{3776AB}
\definecolor{pgblue}{HTML}{336791}
\definecolor{tlsgreen}{HTML}{2E8B57}
\definecolor{netgray}{HTML}{6C757D}
\definecolor{bgsoft}{HTML}{F5F7FA}
\definecolor{bordersoft}{HTML}{C8D1DC}
\definecolor{textdark}{HTML}{1F2937}
\definecolor{accentred}{HTML}{C0392B}
\definecolor{accentyellow}{HTML}{F1C40F}

% Styles génériques réutilisés dans plusieurs diagrammes
\tikzset{
  base/.style={
    rounded corners=4pt,
    draw=bordersoft,
    fill=white,
    text=textdark,
    align=center,
    inner sep=6pt,
    font=\small
  },
  zone/.style={
    base,
    fill=bgsoft,
    draw=netgray,
    dashed,
    inner sep=12pt
  },
  cloud/.style={
    base,
    fill=awsdark!5,
    draw=awsorange,
    line width=0.8pt,
    inner sep=14pt
  },
  service/.style={
    base,
    line width=0.6pt,
    minimum width=2.6cm,
    minimum height=1.1cm,
    drop shadow={shadow xshift=1pt, shadow yshift=-1pt, opacity=0.2}
  },
  legendebox/.style={
    base,
    fill=white,
    font=\footnotesize,
    inner sep=4pt
  },
  fleche/.style={
    -{Latex[length=2.5mm,width=2mm]},
    line width=0.6pt,
    color=netgray
  },
  fleche-bidir/.style={
    {Latex[length=2.5mm,width=2mm]}-{Latex[length=2.5mm,width=2mm]},
    line width=0.6pt,
    color=netgray
  },
  fleche-bold/.style={
    -{Latex[length=3mm,width=2.4mm]},
    line width=1pt
  },
  etiquette/.style={
    font=\scriptsize\sffamily,
    fill=white,
    inner sep=2pt,
    text=textdark
  }
}
 après \documentclass{standalone}
\usepackage[utf8]{inputenc}
\usepackage[T1]{fontenc}
\usepackage[french]{babel}
\usepackage{lmodern}
\usepackage{tikz}
\usepackage{xcolor}

\usetikzlibrary{
  arrows.meta,
  positioning,
  shapes.geometric,
  shapes.symbols,
  shapes.misc,
  fit,
  calc,
  backgrounds,
  decorations.pathreplacing,
  decorations.markings,
  shadows.blur,
  patterns
}

% Palette de couleurs (cohérente sur tout le rapport)
\definecolor{awsorange}{HTML}{FF9900}
\definecolor{awsdark}{HTML}{232F3E}
\definecolor{awsblue}{HTML}{1A476F}
\definecolor{dockerblue}{HTML}{2496ED}
\definecolor{dockerdark}{HTML}{0DB7ED}
\definecolor{nginxgreen}{HTML}{009639}
\definecolor{flaskgray}{HTML}{3776AB}
\definecolor{pgblue}{HTML}{336791}
\definecolor{tlsgreen}{HTML}{2E8B57}
\definecolor{netgray}{HTML}{6C757D}
\definecolor{bgsoft}{HTML}{F5F7FA}
\definecolor{bordersoft}{HTML}{C8D1DC}
\definecolor{textdark}{HTML}{1F2937}
\definecolor{accentred}{HTML}{C0392B}
\definecolor{accentyellow}{HTML}{F1C40F}

% Styles génériques réutilisés dans plusieurs diagrammes
\tikzset{
  base/.style={
    rounded corners=4pt,
    draw=bordersoft,
    fill=white,
    text=textdark,
    align=center,
    inner sep=6pt,
    font=\small
  },
  zone/.style={
    base,
    fill=bgsoft,
    draw=netgray,
    dashed,
    inner sep=12pt
  },
  cloud/.style={
    base,
    fill=awsdark!5,
    draw=awsorange,
    line width=0.8pt,
    inner sep=14pt
  },
  service/.style={
    base,
    line width=0.6pt,
    minimum width=2.6cm,
    minimum height=1.1cm,
    drop shadow={shadow xshift=1pt, shadow yshift=-1pt, opacity=0.2}
  },
  legendebox/.style={
    base,
    fill=white,
    font=\footnotesize,
    inner sep=4pt
  },
  fleche/.style={
    -{Latex[length=2.5mm,width=2mm]},
    line width=0.6pt,
    color=netgray
  },
  fleche-bidir/.style={
    {Latex[length=2.5mm,width=2mm]}-{Latex[length=2.5mm,width=2mm]},
    line width=0.6pt,
    color=netgray
  },
  fleche-bold/.style={
    -{Latex[length=3mm,width=2.4mm]},
    line width=1pt
  },
  etiquette/.style={
    font=\scriptsize\sffamily,
    fill=white,
    inner sep=2pt,
    text=textdark
  }
}
 après \documentclass{standalone}
\usepackage[utf8]{inputenc}
\usepackage[T1]{fontenc}
\usepackage[french]{babel}
\usepackage{lmodern}
\usepackage{tikz}
\usepackage{xcolor}

\usetikzlibrary{
  arrows.meta,
  positioning,
  shapes.geometric,
  shapes.symbols,
  shapes.misc,
  fit,
  calc,
  backgrounds,
  decorations.pathreplacing,
  decorations.markings,
  shadows.blur,
  patterns
}

% Palette de couleurs (cohérente sur tout le rapport)
\definecolor{awsorange}{HTML}{FF9900}
\definecolor{awsdark}{HTML}{232F3E}
\definecolor{awsblue}{HTML}{1A476F}
\definecolor{dockerblue}{HTML}{2496ED}
\definecolor{dockerdark}{HTML}{0DB7ED}
\definecolor{nginxgreen}{HTML}{009639}
\definecolor{flaskgray}{HTML}{3776AB}
\definecolor{pgblue}{HTML}{336791}
\definecolor{tlsgreen}{HTML}{2E8B57}
\definecolor{netgray}{HTML}{6C757D}
\definecolor{bgsoft}{HTML}{F5F7FA}
\definecolor{bordersoft}{HTML}{C8D1DC}
\definecolor{textdark}{HTML}{1F2937}
\definecolor{accentred}{HTML}{C0392B}
\definecolor{accentyellow}{HTML}{F1C40F}

% Styles génériques réutilisés dans plusieurs diagrammes
\tikzset{
  base/.style={
    rounded corners=4pt,
    draw=bordersoft,
    fill=white,
    text=textdark,
    align=center,
    inner sep=6pt,
    font=\small
  },
  zone/.style={
    base,
    fill=bgsoft,
    draw=netgray,
    dashed,
    inner sep=12pt
  },
  cloud/.style={
    base,
    fill=awsdark!5,
    draw=awsorange,
    line width=0.8pt,
    inner sep=14pt
  },
  service/.style={
    base,
    line width=0.6pt,
    minimum width=2.6cm,
    minimum height=1.1cm,
    drop shadow={shadow xshift=1pt, shadow yshift=-1pt, opacity=0.2}
  },
  legendebox/.style={
    base,
    fill=white,
    font=\footnotesize,
    inner sep=4pt
  },
  fleche/.style={
    -{Latex[length=2.5mm,width=2mm]},
    line width=0.6pt,
    color=netgray
  },
  fleche-bidir/.style={
    {Latex[length=2.5mm,width=2mm]}-{Latex[length=2.5mm,width=2mm]},
    line width=0.6pt,
    color=netgray
  },
  fleche-bold/.style={
    -{Latex[length=3mm,width=2.4mm]},
    line width=1pt
  },
  etiquette/.style={
    font=\scriptsize\sffamily,
    fill=white,
    inner sep=2pt,
    text=textdark
  }
}


\begin{document}
\begin{tikzpicture}[
  acteur/.style={
    base, fill=awsblue!10, draw=awsblue, line width=0.6pt,
    minimum width=2.4cm, minimum height=0.8cm,
    font=\small\bfseries
  },
  msg/.style={
    -{Latex[length=2mm]}, line width=0.6pt, color=textdark
  },
  msg-ret/.style={
    -{Latex[length=2mm]}, dashed, line width=0.6pt, color=netgray
  },
  activation/.style={
    fill=awsorange!40, draw=awsorange, line width=0.4pt
  },
  note/.style={
    base, fill=accentyellow!15, draw=accentyellow!80!black,
    font=\scriptsize, align=left, inner sep=3pt
  }
]

  % Positions horizontales des lignes de vie
  \def\xUser{0}
  \def\xBrowser{3.2}
  \def\xNginx{6.4}
  \def\xFlask{9.6}
  \def\xDB{12.8}

  % Acteurs
  \node[acteur] (user)    at (\xUser,    0) {Utilisateur};
  \node[acteur] (browser) at (\xBrowser, 0) {Navigateur};
  \node[acteur] (nginx)   at (\xNginx,   0) {Nginx};
  \node[acteur] (flask)   at (\xFlask,   0) {Flask};
  \node[acteur] (db)      at (\xDB,      0) {PostgreSQL};

  % Lignes de vie (pointillées)
  \def\bottom{-12.5}
  \foreach \x/\name in {\xUser/user, \xBrowser/browser, \xNginx/nginx, \xFlask/flask, \xDB/db} {
    \draw[dashed, color=netgray] (\x, -0.4) -- (\x, \bottom);
  }

  % --- Étape 1 : saisie identifiants ---
  \draw[msg] (\xUser, -1.0) -- node[etiquette, above]{1. Saisit login + password}
    (\xBrowser, -1.0);

  % --- Étape 2 : POST /login ---
  \draw[msg] (\xBrowser, -2.0) -- node[etiquette, above]{2. POST \texttt{/login} (HTTPS)}
    (\xNginx, -2.0);

  % Activation Nginx
  \fill[activation] (\xNginx-0.1, -2.0) rectangle (\xNginx+0.1, -10.0);

  % --- Étape 3 : Terminaison TLS + proxy_pass ---
  \draw[msg] (\xNginx, -2.8) -- node[etiquette, above]{3. proxy\_pass \texttt{http://app:5000}}
    (\xFlask, -2.8);

  % Activation Flask
  \fill[activation] (\xFlask-0.1, -2.8) rectangle (\xFlask+0.1, -9.2);

  % --- Étape 4 : CSRF token check ---
  \draw[msg] (\xFlask, -3.6) to[bend left=90, looseness=4]
    node[etiquette, right, xshift=2mm]{4. Validation CSRF + form}
    (\xFlask, -4.2);

  % --- Étape 5 : SELECT user ---
  \draw[msg] (\xFlask, -5.0) -- node[etiquette, above]{5. \texttt{SELECT * FROM users WHERE username=?}}
    (\xDB, -5.0);

  % Activation DB
  \fill[activation] (\xDB-0.1, -5.0) rectangle (\xDB+0.1, -6.2);

  % --- Étape 6 : retour user record ---
  \draw[msg-ret] (\xDB, -6.0) -- node[etiquette, above, text=netgray]{6. tuple user (id, password\_hash)}
    (\xFlask, -6.0);

  % --- Étape 7 : check_password (bcrypt) ---
  \draw[msg] (\xFlask, -7.0) to[bend left=90, looseness=4]
    node[etiquette, right, xshift=2mm]{7. \texttt{check\_password\_hash()}}
    (\xFlask, -7.6);

  % --- Étape 8 : Set-Cookie session ---
  \draw[msg-ret] (\xFlask, -8.4) -- node[etiquette, above, text=tlsgreen]{8. 302 + \texttt{Set-Cookie: session=...} (HttpOnly, Secure)}
    (\xNginx, -8.4);

  \draw[msg-ret] (\xNginx, -9.2) -- node[etiquette, above, text=tlsgreen]{9. réponse HTTPS}
    (\xBrowser, -9.2);

  \draw[msg-ret] (\xBrowser, -10.0) -- node[etiquette, above]{10. Redirection \texttt{/notes}}
    (\xUser, -10.0);

  % --- Note de sécurité ---
  \node[note, anchor=north west, text width=10cm]
    at (1.5, -11.0)
    {\textbf{Notes de sécurité :}\\
     \(\bullet\) Le mot de passe ne quitte jamais le serveur en clair (HTTPS de bout en bout).\\
     \(\bullet\) Hachage \texttt{Werkzeug PBKDF2-SHA256} avec sel automatique.\\
     \(\bullet\) Cookie de session signé (HMAC) avec \texttt{SECRET\_KEY} et flags \texttt{Secure}, \texttt{HttpOnly}, \texttt{SameSite=Lax}.};

\end{tikzpicture}
\end{document}
